% Mathematical proposal. No claim of a completed R16 experiment.
% Uses the original R15 finite capital economy and its fixed reference policy.
\subsection{A signed occupation account with shared continuation paths}
\label{sec:r16signedbridge}

Fix the finite economy in \eqref{eq:r15finite}--\eqref{eq:r15reward} and
condition on the fitted candidate, the continuation predictor, and their
selection. Let $F_k(x;\xi)$ be the reference continuation payoff along one
future innovation array, starting at the postdecision mean $x$, including
the innovation at $k$, excluding the current reward, and using the original
global grid thereafter. Thus $\E F_k(x;\xi)=S_k^\pi(x)$. The two endpoint
rollouts below use the same array $\xi$. For $r$ independent arrays, average
each array and its exact negative, and denote this average by
$\overline F_k(x)$. The innovation clipping map is odd. Antithetic pairing
preserves the marginal law and hence the expectation of the continuation.

Draw an independent trajectory under the selected candidate and an independent
uniform date $K$. At that date write $\delta=a_K-\pi_K\bm1$ and
$x_a=x_{a_K}(Z_K)$, $x_\pi=x_{\pi_K}(Z_K)$. Define
\begin{align}
 G&=N\{\ell_K(Z_K,a_K)-\ell_K(Z_K,\pi_K)
                 +\overline F_K(x_a)-\overline F_K(x_\pi)\},\label{eq:r16signedG}\\
 C&=N\{\widehat S_K(x_a)-\widehat S_K(x_\pi)
                 -\overline F_K(x_a)+\overline F_K(x_\pi)\}.
 \label{eq:r16signedC}
\end{align}
The difference of the two continuation payoffs is evaluated by differencing
their running production rewards and their terminal dispersion rewards,
using common innovations throughout. In exact arithmetic, $G=M-C$ for every
sample. This cancellation is applied before clipping or constructing a
confidence interval.

\begin{proposition}[Signed Bellman assessment]
\label{prop:r16signedbridge}
Under the original integrability assumptions,
\[
 \E G=\mathcal R_h(a)-\mathcal R_h(\pi),\qquad
 \E C=\mathcal M_h-\{\mathcal R_h(a)-\mathcal R_h(\pi)\}.
\]
Consequently, a fresh simultaneous lower endpoint $L_G\leq\E G$, including
its numerical and clipping allowances, implies
\[
 J(a)-J(m_0)\geq L_G-D_{\rm hold}-E_{a,h}^{\rm all}.
\]
No positive-part replacement of $C$ is made. No separate endpoints for $M$
and $C$ are subtracted to obtain $L_G$.
\end{proposition}

\begin{proof}
Conditional on the candidate history through $K$, antithetic averaging is
unbiased at each endpoint, and their shared innovations affect covariance
only. Hence the conditional expectation of $G/N$ is the reference Bellman
advantage $Q_K^\pi(Z_K,a_K)-Q_K^\pi(Z_K,\pi_K)$. Uniform sampling of $K$
and the finite performance-difference identity \eqref{eq:r15telescoping}
give the first assertion. The second follows from the sample identity
$C=M-G$. The same finite-reference holding deficit and continuous-payoff
transfer used in Theorem~\ref{thm:r15mechanism} give the final assertion.
\end{proof}

The proposition identifies the contribution of the fitted continuation to
the advantage it predicts, while retaining the signed error of that prediction.
A positive $L_G$ verifies the selected candidate's occupation Bellman account.
It does not, by itself, rank the learned critic against a Raw estimator or
attribute a package-level gain uniquely to continuation learning. Those
comparisons require the separately specified common-data experiment.

\subsection{A range that contains no fitted network norm}
\label{app:r16signedrange}

The antithetic difference in \eqref{eq:r16signedG} has a range based only on
the economy and the action tube. Let $\Delta$ be a proved uniform bound on
$\|a_k-\pi_k\bm1\|_d$, and retain the original bound $\varepsilon$ on the
candidate's deviation from the scalar schedule. Let $s_0$ be the maximal
initial dispersion. Put $\gamma=2\chi e^{-\rho T}$, $s_d=\sqrt{d-1}$,
$t_k=kh$, $n_k=N-k-1$, $\tau_k=(N-k)h$, and
\[
 F_k=(1+h\beta)^{n_k},\qquad
 P_k=\beta_\psi\sum_{j=k+1}^{N-1} B_j(1+h\beta)^{j-k-1},
\qquad H_k=s_0+(\kappa+\varepsilon)(t_k+h).
\]
Here $\beta\geq\kappa\|B\|_2$ and
$\beta_\psi\geq\sup_x\|Df(x)^\top\bm1\|_d$ are independently
proved bounds, as in the paired-transfer account. The bound $H_k$ covers
the prefix and the current postdecision bridge. Projection removes the
scalar reference action and the exogenous scalar drift; the centered
candidate action has norm at most $\varepsilon$. A separately proved
state-representation allowance is added to $H_k$ when the implemented
bridge endpoints require it.

Define, for the single future bank of $r$ antithetic pairs,
\begin{align}
 C_{Z,k}&=P_k+\gamma F_k\kappa n_kh
       +\gamma(F_k-1)\left\{H_k+
          \frac{\sigma_I s_d}{\sqrt d}(\sqrt{t_k}+\sqrt{\tau_k})\right\},\notag\\
 A_{Z,k}&=\gamma(F_k-1)\frac{\sigma_I}{\sqrt d}
                         (\sqrt{t_k}+\sqrt{\tau_k/r}).
 \label{eq:r16residualrange}
\end{align}
Let $L_{\ell,k}$ be a proved bound on the normalized action-gradient
norm of $\ell_k/h$ throughout the original action box. For example,
\[
 L_{\ell,k}=\max_{m\in\{l_k,u_k\}}
       |(A_k/h)(m^{-1}-\zeta m)-B_k/h|
\]
is valid when $[l_k,u_k]$ bounds every coordinate of the action and its
line segment to the reference. The derivative is decreasing in its own
consumption coordinate and in the cross-sectional mean, so the two
diagonal endpoints enclose every component.
Set
\begin{align}
 C_{G,k}&=T\Delta\left[L_{\ell,k}
       +\gamma\left(H_k+\frac{\sigma_I s_d\sqrt{t_k}}{\sqrt d}\right)
       +C_{Z,k}\right],\notag\\
 A_{G,k}&=T\Delta\left[\frac{\gamma\sigma_I\sqrt{t_k}}{\sqrt d}
                                      +A_{Z,k}\right],\qquad
 C_G=\max_k C_{G,k},\quad A_G=\max_k A_{G,k}.
 \label{eq:r16signedrange}
\end{align}

\begin{proposition}[Antithetic paired-continuation range]
\label{prop:r16signedrange}
For the exact-real statistic in \eqref{eq:r16signedG},
\[
 \Pr\{|G|>C_G+A_G v\}\leq 2e^{-v^2/2},\qquad v\geq0.
\]
For a fixed $v_*>0$, valid clipping and expectation-tail allowances are
\[
 b_G=C_G+A_Gv_*,\qquad
 \tau_G=2A_Ge^{-v_*^2/2}/v_*.
\]
The same formulas may be evaluated conditionally on a prespecified date
or initial-profile stratum before averaging its protected endpoint.
They require no maximum of an observed fitted derivative or simulated
outcome. All executed arithmetic errors are added separately.
\end{proposition}

\begin{proof}
For each fixed innovation array, the mean-value identity gives
\[
 \overline F_k(x_a)-\overline F_k(x_\pi)
       =-h\int_0^1\langle\delta,
               \overline q_k(x_\pi-uh\delta)\rangle_d\,du,
\]
where $\overline q_k=dD\overline F_k$. This is an identity used for the
proof; computing the statistic requires no derivatives.

The adjoint recursion in \eqref{eq:r15adjoint} gives
$\|J\|_2\leq F_k$, $\|J-I\|_2\leq F_k-1$, and production-adjoint
norm at most $P_k$. For the terminal derivative, subtract
$q_0(x)=-\gamma\Pi x$. The exact residual consists of the production
adjoint and
\[
 -\gamma\left\{(J^\top-I)\Pi x
       +J^\top\sum_{j>k}h\Pi f(Z_j)
       +(J^\top-I)\Pi\xi+\Pi\xi\right\}.
\]
The average of the final term is exactly zero in every antithetic pair.
The other terms are bounded by $P_k$, $\gamma F_k\kappa n_kh$, and
$\gamma(F_k-1)(\|\Pi x\|_d+\|\Pi\xi\|_d)$.

The normalized norm of the projected prefix innovation sum has mean
at most $s_d$ and Lipschitz constant one as a function of its standard
Gaussian inputs. The average of the corresponding future norms from
$r$ independent base paths has mean at most $s_d$ and Lipschitz constant
$r^{-1/2}$. Odd coordinatewise clipping is one-Lipschitz. Gaussian
concentration and a union bound therefore give a common event of failure
probability at most $2e^{-v^2/2}$ on which, uniformly in $u$,
\[
 \|\Pi x_u\|_d\leq H_k+
               \frac{\sigma_I\sqrt{t_k}}{\sqrt d}(s_d+v),
 \quad
 \|\overline q_k(x_u)+\gamma\Pi x_u\|_d
                                      \leq C_{Z,k}+A_{Z,k}v.
\]
The future event is uniform in the endpoint and the bridge parameter:
only the common innovation norms and global drift/adjoint bounds enter.
No union over bridge points or integration nodes is required. For $k=0$
the prefix contribution is zero.

The known terminal derivative integrates exactly:
\[
 -h\int_0^1\langle\delta,-\gamma\Pi x_u\rangle_d\,du
       =h\gamma\langle\delta,\Pi(x_a+x_\pi)/2\rangle_d.
\]
Combining its bound, the residual bound, and
$N|\ell_k(z,a)-\ell_k(z,\pi)|\leq T\Delta L_{\ell,k}$
proves the displayed tail inequality. Integrating that inequality above
$b_G$ and using
$\int_{v_*}^{\infty}e^{-v^2/2}dv\leq e^{-v_*^2/2}/v_*$ proves
the expectation allowance. The bound concerns the exact-real statistic;
it does not remove roundoff, protected-function, or state-representation
allowances from an implementation.
\end{proof}

\paragraph*{The asymmetric stage range used in the implementation.}
The implemented account further avoids symmetrizing the known current
reward. Define
\[
 f_k(m)=A_k\log m-B_k m-\zeta A_km^2/2,
\]
\[
 L_k^{\rm stage}=N\min\{f_k(l_k)-f_k(\pi_k),
                              f_k(u_k)-f_k(\pi_k)\},
\]
and
\[
 U_k^{\rm stage}=N\left|A_k(\pi_k^{-1}-\zeta\pi_k)-B_k\right|\Delta_k,
 \quad \Delta_k\geq\sup\|a_k-\pi_k\bm1\|_d.
\]
These are respectively lower and upper bounds on the current stage
advantage. For the lower bound, the chord of $\log$ on $[l_k,u_k]$
is a lower bound, and its coordinate average depends only on the
cross-sectional mean. Subtracting the linear term and the mean quadratic
gives a concave function of that mean, whose minimum is at an endpoint.
For the upper bound, apply the concave reward's tangent support at the
uniform reference action and Cauchy--Schwarz.

Put
\[
 C_{q,k}=C_{Z,k}+\gamma\left(H_k+\sigma_I s_d\sqrt{t_k/d}\right),
 \qquad A_{q,k}=A_{Z,k}+\gamma\sigma_I\sqrt{t_k/d}.
\]
On the same two-event intersection,
\[
 L_k^{\rm stage}-T\Delta_k(C_{q,k}+A_{q,k}v)
       \leq G\leq
 U_k^{\rm stage}+T\Delta_k(C_{q,k}+A_{q,k}v).
\]
The minimum lower endpoint and maximum upper endpoint at the fixed $v_*$
therefore give a valid asymmetric clipping interval. Centering that interval
at an exact binary64 representative and enclosing both distances from its
center changes neither the estimand nor the event. The absolute expectation
allowance is
$2\max_k(T\Delta_k A_{q,k})e^{-v_*^2/2}/v_*$.
Thus the tighter date-specific prefix and asymmetric stage ranges retain
the same proof and spend no additional probability budget.
